% Apendice A (revisao pos-EE). Gerado por 02_audit/build12_appendix.py; tabulars verbatim.
\begin{table}[htbp]\centering
\caption{Income gradient in one-year referral by birth cohort and postcode income quartile}
\label{tab:E2}
\label{tab:P2grad}
\footnotesize
\par\medskip\textbf{Panel A. Referral incidence and restricted mean time}\par\smallskip
\fitbox{%
\begin{tabular}{lrrrrrr}
\toprule
Cohort & ref Q1 & ref Q2 & ref Q3 & ref Q4 & RMST Q1 & RMST Q4 \\
\midrule
2010 & 0.530 & 0.520 & 0.566 & 0.652 & 222 & 199 \\
2011 & 0.505 & 0.483 & 0.530 & 0.613 & 252 & 234 \\
2012 & 0.607 & 0.596 & 0.625 & 0.701 & 240 & 222 \\
2013 & 0.722 & 0.722 & 0.760 & 0.815 & 212 & 182 \\
2014 & 0.734 & 0.729 & 0.757 & 0.828 & 222 & 199 \\
2015 & 0.790 & 0.782 & 0.812 & 0.876 & 199 & 179 \\
2016 & 0.855 & 0.853 & 0.877 & 0.923 & 184 & 167 \\
\bottomrule
\end{tabular}}
\par\medskip\textbf{Panel B. Full series and the gap between quartiles}\par\smallskip
\fitbox{%
\begin{tabular}{lrrr}
\toprule
Cohort & Q1 & Q4 & gap (p.p.) \\
\midrule
2010 & 0.530 & 0.652 & +12.2 \\
2011 & 0.505 & 0.613 & +10.8 \\
2012 & 0.607 & 0.701 & +9.4 \\
2013 & 0.722 & 0.815 & +9.3 \\
2014 & 0.734 & 0.828 & +9.4 \\
2015 & 0.790 & 0.876 & +8.6 \\
2016 & 0.855 & 0.923 & +6.8 \\
\bottomrule
\end{tabular}}
\tabnote{\emph{Source:} SME childcare request record, 2010--2018 applications, with referral and enrollment dates through November 2019. Income: mean household income of the postcode's census tract (2010 Census), not household income. \emph{Sample:} duration sample of Table~\ref{tab:E1cohort}. \emph{Variables:} cumulative incidence of referral within 365 days by postcode income quartile (Q1 lowest); RMST in days; Q4$-$Q1 gap in percentage points. \emph{Exercise:} Aalen--Johansen estimator with aging out competing; descriptive.}
\end{table}
